\chapter{Oscillator Wavefunctions, Spin-One-Half, and Particle Exchange}

\section{Oscillators Everywhere}

Any system with a stable equilibrium can oscillate when it is displaced from that equilibrium. It may have one frequency or many. A violin string, for example, has a fundamental mode and higher harmonics. Plucking it excites some combination of those modes.

The more important examples here are fields. The equilibrium configuration of the electromagnetic field is the vacuum: no electric field, no magnetic field, and no electromagnetic oscillation. Send radiation into a cavity or a waveguide and the field can oscillate in its allowed modes, each with its own frequency. In the quantum theory those modes describe collections of photons.

The same idea applies to a crystal. It is misleading to picture every atom as an independent oscillator. The atoms exert forces on one another, so a displacement of one atom affects its neighbors and a wave travels through the crystal. The normal modes of that coupled motion are harmonic oscillators, and their quanta are called phonons. Photons and phonons are different physical objects, but both lead back to the same oscillator algebra.

\section{The Oscillator Algebra}

For the oscillator discussion, set the mass and the reduced Planck constant equal to one:
\[
m=1,\qquad \hbar=1.
\]
The Hamiltonian is
\begin{equation}
H=\frac{p^{2}}{2}+\frac{\omega^{2}x^{2}}{2}.
\end{equation}
The parameter \(\omega\) fixes the curvature of the potential and is also the oscillation frequency.

Define the raising and lowering operators by
\begin{equation}
A_{\pm}=\frac{p\pm i\omega x}{\sqrt{2\omega}}.
\end{equation}
The normalization is chosen so that
\begin{equation}
[A_{-},A_{+}]=1.
\end{equation}
The number operator is
\begin{equation}
N=A_{+}A_{-}.
\end{equation}
Together, these equations imply that \(A_{+}\) raises an eigenvalue of \(N\) by one and \(A_{-}\) lowers it by one. The resulting spectrum is
\begin{equation}
n=0,1,2,3,\ldots .
\end{equation}

Multiplying the two ladder operators also displays the quantum correction to the classical energy:
\begin{align}
N
&=\frac{1}{2\omega}(p+i\omega x)(p-i\omega x)\\
&=\frac{1}{2\omega}\left(p^{2}+\omega^{2}x^{2}
 +i\omega[x,p]\right)\\
&=\frac{1}{2\omega}\left(p^{2}+\omega^{2}x^{2}-\omega\right),
\end{align}
where \([x,p]=i\). Consequently,
\begin{equation}
H=\omega\left(N+\frac12\right).
\end{equation}
If the additive constant is irrelevant to the problem, it is convenient to use the shifted Hamiltonian
\begin{equation}
H_{\mathrm{shifted}}=H-\frac{\omega}{2}=\omega N.
\end{equation}
The term \(\omega/2\) is the oscillator's zero-point, or ground-state, energy. For a field mode the analogous term is called vacuum energy. It reflects the uncertainty principle: no state can have both \(x\) and \(p\) sharply equal to zero, so the two positive quadratic terms in \(H\) cannot both have zero expectation value.

\begin{classroomqa}
\lecturetimestamp{00:09:46}
\audiencequestion{Why is zero included in the spectrum of the number operator?}
\lecturerresponse{It can be demonstrated explicitly by constructing the ground-state wavefunction. That construction will exhibit a state with \(N=0\).}
\end{classroomqa}

To begin that construction, replace state vectors by wavefunctions of position. Energy eigenstates become
\[
\psi_{0}(x),\ \psi_{1}(x),\ \psi_{2}(x),\ldots,
\]
and the basic operators act as
\begin{equation}
X\longmapsto x,\qquad
P\longmapsto -i\frac{d}{dx}.
\end{equation}
With \(\hbar\) restored, the second replacement would be
\(P\mapsto-i\hbar\,d/dx\).

Before solving for \(\psi_{0}\), there is a faster abstract answer to the question. On the discrete ladder of normalizable oscillator eigenstates, \(N=A_{+}A_{-}\) is nonnegative, so repeated lowering in unit steps cannot continue indefinitely. At the endpoint---whose normalizable wavefunction will be constructed below---there must be a state, denoted \(\lvert0\rangle\), for which
\begin{equation}
A_{-}\lvert0\rangle=\mathbf 0.
\end{equation}
Otherwise \(A_{-}\lvert0\rangle\) would be another state one step lower. It follows immediately that
\begin{equation}
N\lvert0\rangle
=A_{+}A_{-}\lvert0\rangle
=\mathbf 0
=0\,\lvert0\rangle.
\end{equation}
Thus the bottom state has number eigenvalue zero.

\begin{classroomqa}
\lecturetimestamp{00:14:57}
\audiencequestion{If zero is the eigenvalue, should the right-hand side contain a zero eigenvector?}
\lecturerresponse{The ket \(\lvert0\rangle\) is a normalized, nonzero vector whose label is zero. In \(N\lvert0\rangle=0\,\lvert0\rangle\), the right-hand side is the zero vector. The ground-state ket and the zero vector are different objects.}
\end{classroomqa}

\section{The Schr\"odinger Equation and Physical Solutions}

Return now to the position representation. The time-independent Schr\"odinger equation is
\begin{equation}
\left[
-\frac12\frac{d^{2}}{dx^{2}}
+\frac12\omega^{2}x^{2}
\right]\psi(x)=E\psi(x).
\end{equation}
As a second-order differential equation, it has formal solutions for arbitrary values of \(E\). Most of those formal solutions are not physical oscillator states. They grow exponentially at large \(\lvert x\rvert\), so their norm diverges.

For a bound state, the wavefunction must be square integrable:
\begin{equation}
\int_{-\infty}^{\infty}\psi^{*}(x)\psi(x)\,dx<\infty.
\end{equation}
It can then be normalized so that the integral equals one. Only special energy values produce such solutions. Position and momentum states on an infinite line are exceptional in this respect, but exponentially growing oscillator solutions are simply excluded from the state space.

\section{The Gaussian Ground State}

The second-order equation can be solved, but the ladder algebra gives a simpler first-order equation:
\begin{equation}
A_{-}\psi_{0}(x)=0.
\end{equation}
The overall factor in \(A_{-}\) is irrelevant when the right-hand side is zero, so this is equivalent to
\begin{equation}
(p-i\omega x)\psi_{0}(x)=0.
\end{equation}
Using \(p=-i\,d/dx\) gives
\begin{equation}
\left(\frac{d}{dx}+\omega x\right)\psi_{0}(x)=0.
\end{equation}

\begin{classroomqa}
\lecturetimestamp{00:22:32}
\audiencequestion{Could this be solved by separation of variables?}
\lecturerresponse{There is only one independent variable, so there are not two variables to separate. Written as \(\psi_{0}'/\psi_{0}=-\omega x\), the equation has the form of a logarithmic derivative. That suggests writing the wavefunction as an exponential.}
\end{classroomqa}

A useful way to solve it is to write
\begin{equation}
\psi_{0}(x)=e^{f(x)}.
\end{equation}
Then
\[
\psi_{0}'(x)=f'(x)e^{f(x)},
\]
and the common exponential factor cancels. The equation reduces to
\begin{equation}
\frac{df}{dx}+\omega x=0.
\end{equation}

\begin{classroomqa}
\lecturetimestamp{00:26:28}
\audiencequestion{Should there be a minus sign in \(f(x)\)?}
\lecturerresponse{Yes. Since \(df/dx=-\omega x\), integration gives a negative quadratic term.}
\end{classroomqa}

Therefore
\begin{equation}
f(x)=-\frac12\omega x^{2}+C,
\end{equation}
and
\begin{equation}
\psi_{0}(x)\propto e^{-\omega x^{2}/2}.
\end{equation}
The multiplicative normalization constant is left unspecified. The important point is the Gaussian shape: it is smooth, concentrated near the origin, and rapidly decreasing at large \(\lvert x\rvert\). It is plainly square integrable.

There is no need to substitute the answer back into the second-order equation to discover its energy. The algebra already gives
\begin{equation}
E_{0}=\frac{\omega}{2}.
\end{equation}

\section{Excited States, Nodes, and Classical Motion}

The first excited state follows by raising the ground state:
\begin{equation}
\psi_{1}(x)\propto A_{+}\psi_{0}(x)
\propto(p+i\omega x)e^{-\omega x^{2}/2}.
\end{equation}
The ground-state equation says
\[
p\psi_{0}=i\omega x\,\psi_{0}.
\]
Changing the sign in the ladder operator turns the cancellation into addition:
\begin{equation}
(p+i\omega x)\psi_{0}
=2i\omega x\,\psi_{0}.
\end{equation}
Up to an overall constant,
\begin{equation}
\psi_{1}(x)\propto x\,e^{-\omega x^{2}/2}.
\end{equation}

This wavefunction is odd:
\begin{equation}
\psi_{1}(-x)=-\psi_{1}(x).
\end{equation}
Its probability density is even, but it has a node at the origin,
\begin{equation}
\lvert\psi_{1}(0)\rvert^{2}=0.
\end{equation}
Thus the first-excited-state probability density has a node at \(x=0\).

Repeated action of \(A_{+}\) gives higher-degree polynomials multiplying the same Gaussian envelope.\footnote{The lecture states the higher-degree polynomial factors, the additional node at each step, and alternating symmetry from \(00{:}34{:}25\) to \(00{:}35{:}49\). The Gaussian-times-polynomial form is corroborated by Leonard Susskind and Art Friedman, \emph{Quantum Mechanics: The Theoretical Minimum} (2014), in its discussion of oscillator eigenfunctions.} Up to normalization,
\begin{equation}
\psi_{n}(x)=P_{n}(x)e^{-\omega x^{2}/2}.
\end{equation}
The parity alternates between even and odd, and each step adds a node. As \(n\) increases, the wavefunction spreads farther from the origin and develops more rapid oscillations.

The two features have direct physical meanings. Support farther from the origin corresponds to larger potential energy because the potential grows as \(x^{2}\). Rapid oscillation in position space signals larger momentum. A highly excited oscillator is likely to reach farther from the center; when it passes through the central region, it carries more momentum and its wavefunction oscillates more rapidly there.

\begin{classroomqa}
\lecturetimestamp{00:37:18}
\audiencequestion{Does the harmonic oscillator fail to approach classical behavior when \(n\) becomes large?}
\lecturerresponse{Classical motion is not represented by one stationary energy eigenstate. It appears in a wave packet containing many energy levels. A displaced Gaussian, for example, evolves back and forth like a classical oscillator, and the relative uncertainty becomes smaller as the swing grows larger.}
\end{classroomqa}

Start with the ground-state Gaussian and displace it from the origin. The displaced packet is no longer an energy eigenstate; it is a superposition of many of them. Under the time-dependent Schr\"odinger equation, its components acquire different phases. The packet swings from one side of the potential to the other, developing rapid phase variation near the center and becoming smooth again near a turning point. The larger the excursion compared with the packet's width, the more nearly classical the motion appears.

\begin{classroomqa}
\lecturetimestamp{00:39:29}
\audiencequestion{Should the second excited-state wavefunction pass through zero at the origin?}
\lecturerresponse{No. Its polynomial factor is not simply \(x^{2}\); it also contains a constant term, so the wavefunction need not vanish at \(x=0\).}
\end{classroomqa}

The ladder argument also assumes that \(A_{+}\) and \(A_{-}\) are Hermitian conjugates on the square-integrable oscillator domain, with the boundary behavior needed for the adjoint relation. Exponentially growing formal solutions do not belong to that domain.

\begin{classroomqa}
\lecturetimestamp{00:40:45}
\audiencequestion{Is there a general way to recognize when a Hamiltonian can be solved by these algebraic methods?}
\lecturerresponse{There is no general rule of that kind. Exactly solvable examples are rare. The harmonic oscillator and the ideal nonrelativistic hydrogen atom are fortunate cases in which algebraic methods work.}
\end{classroomqa}

\section{Spin One-Half as Angular Momentum}

Spin is angular momentum carried by a particle independently of its center-of-mass orbital motion. A particle at rest has no orbital angular momentum \(\mathbf R\times\mathbf P\), but it may still have spin. One can try to picture a tiny spinning object, or treat spin as an intrinsic quantum property. The decisive statement is mathematical: its components transform under spatial rotations and obey the angular-momentum commutation relations.

Restore \(\hbar\) from this point onward. Angular-momentum components satisfy
\begin{equation}
[L_{z},L_{x}]=i\hbar L_{y},
\end{equation}
with the other two relations obtained by cyclic permutation. The dimensionless Pauli matrices are
\begin{equation}
\sigma_{z}=
\begin{pmatrix}
1&0\\
0&-1
\end{pmatrix},
\qquad
\sigma_{x}=
\begin{pmatrix}
0&1\\
1&0
\end{pmatrix},
\qquad
\sigma_{y}=
\begin{pmatrix}
0&-i\\
i&0
\end{pmatrix}.
\end{equation}
This basis makes \(\sigma_{z}\) diagonal, with eigenvalues \(+1\) and \(-1\).

\begin{classroomqa}
\lecturetimestamp{00:47:35}
\audiencequestion{Is there a minus sign in \(\sigma_x\)?}
\lecturerresponse{No. The two off-diagonal entries of \(\sigma_x\) are both positive.}
\end{classroomqa}

Multiplication gives\footnote{The entry-by-entry products are partly visible in \texttt{lecture\_04\_frame\_01.png} at \(00{:}49{:}35\); the completed relations follow the calculation spoken immediately around that frame.}
\begin{align}
\sigma_{z}\sigma_{x}
&=
\begin{pmatrix}
0&1\\
-1&0
\end{pmatrix}
=i\sigma_{y},\\
\sigma_{x}\sigma_{z}
&=
\begin{pmatrix}
0&-1\\
1&0
\end{pmatrix}
=-i\sigma_{y}.
\end{align}
Hence
\begin{equation}
[\sigma_{z},\sigma_{x}]=2i\sigma_{y}.
\end{equation}
The extra factor of two is removed by defining the physical spin operators as\footnote{The board calculation suppresses \(\hbar\). It is restored here so that \(S_i\) has the units of angular momentum.}
\begin{equation}
S_{i}=\frac{\hbar}{2}\sigma_{i}.
\end{equation}
Then
\begin{equation}
[S_{z},S_{x}]=i\hbar S_{y},
\end{equation}
and similarly for cyclic permutations. Half the Pauli matrices, with the factor of \(\hbar\) restored, therefore form an angular-momentum system.

Setting \(\hbar=1\) gives the form shown below.\footnote{The source frame is \texttt{lecture\_04\_frame\_02.png}.}
\begin{figure}[htbp]
\centering
\includegraphics[width=0.9\linewidth]{lecture_04_frame_02.png}
\caption{In units with \(\hbar=1\), half the Pauli matrices are the spin-one-half angular-momentum operators and satisfy \([S_z,S_x]=iS_y\).}
\lectureframe{00:55:41}
\end{figure}

The two \(S_z\) eigenvalue equations are
\begin{equation}
S_z\lvert\pm\rangle
=\pm\frac{\hbar}{2}\lvert\pm\rangle.
\end{equation}
Equivalently, the dimensionless magnetic quantum number is
\begin{equation}
m_{s}=\pm\frac12,
\qquad
S_{z}\lvert m_{s}\rangle
=m_{s}\hbar\,\lvert m_{s}\rangle.
\end{equation}
Angular-momentum raising and lowering operators shift \(m\) by integer steps, but this does not require every ladder to pass through zero. Integer multiplets contain integer \(m\)-values, while half-integer multiplets contain half-integer values. For spin one-half the ladder has only the two members \(m_s=+1/2\) and \(m_s=-1/2\).

These two states form the spin-one-half doublet. Electrons, muons, neutrinos, and quarks are among the particles described by such half-integer spin multiplets.

\section{Orbital and Spin Angular Momentum}

There are now two kinds of angular momentum. Orbital angular momentum describes motion relative to an origin,
\begin{equation}
\mathbf L=\mathbf R\times\mathbf P,
\end{equation}
while \(\mathbf S\) is intrinsic spin. Their vector sum is the total angular momentum,
\begin{equation}
\mathbf J=\mathbf L+\mathbf S.
\end{equation}

\begin{classroomqa}
\lecturetimestamp{00:56:49}
\audiencequestion{If a system has total angular momentum \(\mathbf J\), how is its spin \(\mathbf S\) identified?}
\lecturerresponse{The spin and orbital quantum numbers are normally supplied separately. If \(L=0\), all the angular momentum is spin. Otherwise one needs the quantum rules for adding angular momenta, which are postponed here.}
\end{classroomqa}

The addition rules apply equally to orbital angular momentum plus spin, to two spins, or to two orbital angular momenta. The combined system has its own allowed values of total angular momentum, but those rules are not needed yet.

\begin{classroomqa}
\lecturetimestamp{00:58:18}
\audiencequestion{What experimental evidence led to the introduction of spin?}
\lecturerresponse{Spectroscopy supplied evidence, but the periodic table gives a particularly transparent motivation: its shell structure requires electrons to possess an additional two-valued property.}
\end{classroomqa}

\section{Hydrogenic Degeneracy and Shell Counting}

For orbital angular momentum, the \(L^{2}\) eigenvalue equation is
\begin{equation}
L^{2}\lvert\ell,m\rangle
=\hbar^{2}\ell(\ell+1)\lvert\ell,m\rangle,
\end{equation}
and a fixed value of \(\ell\) has
\begin{equation}
2\ell+1
\end{equation}
magnetic states, corresponding to \(m=-\ell,-\ell+1,\ldots,\ell\).

The ideal nonrelativistic Coulomb problem has an additional degeneracy. Its bound-state energies are negative and approach zero from below, but only the alignment of levels matters for the counting. The first \(\ell=1\) level aligns with the next \(\ell=0\) level; the first \(\ell=2\) level joins the next principal shell; and the pattern continues.\footnote{\texttt{lecture\_04\_frame\_03.png}, \(01{:}02{:}01\); the completed principal-shell alignment is stated from \(01{:}02{:}33\) through \(01{:}04{:}10\).}

\begin{figure}[htbp]
\centering
\begin{tikzpicture}[x=1.55cm,y=1.0cm]
\draw[->] (-0.65,0) -- (-0.65,4.35);
\draw[->] (-0.85,0) -- (4.75,0);
\node at (-0.95,4.08) {\(E\)};
\node at (4.55,-0.28) {\(\ell\)};
\node at (0,-0.30) {\(0\)};
\node at (2,-0.30) {\(1\)};
\node at (4,-0.30) {\(2\)};

\draw (-0.20,0.75) -- (0.20,0.75);

\draw (-0.20,2.00) -- (0.20,2.00);
\foreach \dx in {-0.18,0,0.18}
  \draw (2+\dx-0.06,2.00) -- (2+\dx+0.06,2.00);

\draw (-0.20,3.15) -- (0.20,3.15);
\foreach \dx in {-0.18,0,0.18}
  \draw (2+\dx-0.06,3.15) -- (2+\dx+0.06,3.15);
\foreach \dx in {-0.32,-0.16,0,0.16,0.32}
  \draw (4+\dx-0.05,3.15) -- (4+\dx+0.05,3.15);

\node[anchor=west] at (0.35,0.75) {\(n=1\)};
\node[anchor=west] at (2.35,2.00) {\(n=2\)};
\node[anchor=west] at (4.40,3.15) {\(n=3\)};
\end{tikzpicture}
\caption{The first three ideal Coulomb shells combine orbital multiplets of equal principal energy: one state in the first shell, \(1+3\) in the second, and \(1+3+5\) in the third.}
\lectureframe{01:02:01}
\end{figure}

In the principal shell \(n\), the allowed orbital values are\footnote{This compact \(\ell\)-range and the sum below reconstruct the level alignments and the counts \(1,4,9,16,\ldots\) stated from \(01{:}02{:}33\) to \(01{:}04{:}10\). The beginning of the spectrum drawing appears in \texttt{lecture\_04\_frame\_03.png} at \(01{:}02{:}01\).}
\[
\ell=0,1,\ldots,n-1.
\]
The number of spatial states is therefore
\begin{equation}
\sum_{\ell=0}^{n-1}(2\ell+1)=n^{2},
\end{equation}
giving the sequence
\begin{equation}
1,\ 4,\ 9,\ 16,\ldots .
\end{equation}
This is orbital counting only. Each spatial state permits two spin values, so the corresponding idealized totals are
\begin{equation}
2,\ 8,\ 18,\ 32,\ldots=2n^{2}.
\end{equation}

\section{Exclusion and the Two-Valued Electron Property}

Begin with helium in an independent-electron approximation. Its nucleus has twice the proton charge, which rescales the hydrogenic levels and draws the electrons inward. Ignoring electron--electron interactions for the moment, two electrons can occupy the lowest spatial orbital and give a reasonable first approximation to the helium ground state.

Why not put a third electron there as well? That would lower its orbital energy, but it does not happen. The restriction is the Pauli exclusion principle: no two electrons may occupy the same complete one-particle state. Calling this an ordinary repulsive force would be misleading. It is a restriction on the allowed quantum state.

There is an apparent contradiction. The two helium electrons do share the same spatial orbital. The resolution is that a complete electron state contains more than the orbital wavefunction. It also contains a two-valued degree of freedom. Call the values up and down. One electron in the lowest orbital can be spin up and the other spin down; their spatial states agree, but their full quantum states do not.

At this point the counting starts forward toward lithium and then backs up. The backtrack exposes why both exclusion and spin are necessary. Imagine an exclusion rule with no spin at all. Only one electron could occupy the lowest spatial state. A second electron would have to enter one of the four spatial states in the next principal shell, producing several alternatives where the observed helium ground state is unique. That spinless version of exclusion is wrong.

Restore the two spin values. Two electrons fill the lowest orbital with opposite spin. A third electron, as in neutral lithium, must enter the next principal shell. That shell has one \(\ell=0\) spatial state and three \(\ell=1\) spatial states. Before spin is counted, there are four possibilities.

\begin{classroomqa}
\lecturetimestamp{01:14:58}
\audiencequestion{How do the \(s\)- and \(p\)-states give eight possibilities in the first excited shell?}
\lecturerresponse{There is one spatial state with \(\ell=0\) and three with \(\ell=1\), giving four spatial states. Each admits two spin values, so the total is \(2(1+3)=8\).}
\end{classroomqa}

This is the simple counting behind the eight entries in the second row of the periodic table. It is deliberately an idealization: the later discussion will make clear that electron--electron interactions soon spoil the independent hydrogenic-shell picture.

The two-valued property can also be tested as angular momentum. The electron's spin is accompanied by a magnetic moment, so a magnetic field gives its two spin projections different energies. A level that was degenerate in zero field splits into two, just as expected for a small magnetic moment aligned or anti-aligned with the field. The splitting tests whether the added two-valued degree of freedom behaves as angular momentum rather than as an arbitrary binary label.

\section{Photons, Electrons, and the Limits of the Shell Picture}

Particles do not all obey the same occupancy rule. Many photons may occupy the same electromagnetic mode; a classical light wave may be understood as a state containing a large number of photons in that mode. Electrons behave differently: exclusion prevents them from piling into one complete quantum state. This contrast motivates the distinction between the two particle statistics.

\begin{classroomqa}
\lecturetimestamp{01:20:28}
\audiencequestion{What happens for a free stream of electrons rather than electrons bound in atoms?}
\lecturerresponse{The exclusion rule still applies: the electrons cannot all occupy one quantum state. A laser can place many photons in the same mode, but an electron stream cannot be prepared in that way.}
\end{classroomqa}

The restriction is a property of electrons, not a peculiarity of atomic binding.

\begin{classroomqa}
\lecturetimestamp{01:21:12}
\audiencequestion{Is orbital angular momentum specifically rotation about a nucleus?}
\lecturerresponse{Angular momentum is defined relative to a chosen origin. The nucleus is the natural origin in an atom, while a free particle has angular momentum \(\mathbf R\times\mathbf P\) relative to whatever origin is selected.}
\end{classroomqa}

Every discussion of rotation chooses a point about which the rotation is made. An atom supplies a natural center because its nucleus is the center of the approximately central potential. In free space no origin is preferred, although angular momentum can still be calculated about any chosen one.

\begin{classroomqa}
\lecturetimestamp{01:22:34}
\audiencequestion{Why does the periodic table contain repeated rows of eight and then repeated rows of eighteen elements?}
\lecturerresponse{The simple independent-electron shell picture does not explain the full pattern. It breaks down quickly because it neglects the interactions among the electrons.}
\end{classroomqa}

For a large atom, the outer electrons interact with a substantial inner electron cloud as well as with the nucleus. Treating every electron as if it moved independently in an unscreened Coulomb field is no longer sensible.

\begin{classroomqa}
\lecturetimestamp{01:23:33}
\audiencequestion{Why do atoms not become dramatically larger as more electrons are added?}
\lecturerresponse{Adding electrons also adds protons to the nucleus. The increased nuclear charge pulls the inner electrons inward, while those inner electrons screen the nucleus. In this simple picture, an outer valence electron then moves in an effective charge of roughly one.}
\end{classroomqa}

The increased nuclear charge contracts the inner cloud. For an atom with one outer valence electron, the nucleus and the inner electrons nearly cancel at large distance, leaving the valence electron in a field roughly like that of a unit positive charge. Atomic sizes therefore do not simply grow in proportion to the number of electrons.

\begin{classroomqa}
\lecturetimestamp{01:24:49}
\audiencequestion{Is electrical repulsion the reason neighboring atoms remain separated?}
\lecturerresponse{Not by itself. Pauli exclusion supplies an important effective repulsion when electron clouds overlap, and for many atoms it is the dominant reason they resist further overlap.}
\end{classroomqa}

Pauli repulsion is an effective consequence of exclusion, not an extra force law.

\begin{classroomqa}
\lecturetimestamp{01:25:23}
\audiencequestion{Is the exclusion principle a separate postulate?}
\lecturerresponse{In nonrelativistic quantum mechanics it is imposed as a postulate. In relativistic quantum theory, the spin--statistics connection gives it a deeper basis.}
\end{classroomqa}

\section{Identical Particles in Classical and Quantum Mechanics}

Before formulating that postulate, consider how identical particles are counted classically. Put two particles of the same kind into a set of boxes. If the particles secretly carry names---Harry and Sally---then Harry in the first box and Sally in the second appears different from Sally in the first and Harry in the second.

Classical mechanics permits that picture even when the labels are too faint to detect. One can imagine an arbitrarily weak mark on each particle and still regard their trajectories as distinct. In classical statistical mechanics one usually chooses not to count permutations of identical particles as new configurations. For \(N\) particles this convention produces the familiar factor \(1/N!\).

Quantum mechanics is stricter. Identical particles do not carry hidden individual labels that can be followed in principle. Two electrons are interchangeable as electrons. Exchanging an electron with a photon would produce a physically different arrangement because they are different species, but exchanging the two electron slots cannot reveal which electron was in fact which.

A one-particle state has a wavefunction
\[
\psi(x)=\langle x\vert\psi\rangle.
\]
For two particles the state is a function of two position arguments,
\begin{equation}
\psi(x_{1},x_{2}).
\end{equation}
The subscripts label the two slots in the description, not two spatial directions. Its absolute square gives the joint position probability density.

\section{The Swap Operator and the Statistics Postulate}

Interchange the two particle slots:
\begin{equation}
\psi(x_{1},x_{2})\longmapsto\psi(x_{2},x_{1}).
\end{equation}
Following the notation used in this discussion, call the swap operator \(S\). Here \(S\) denotes exchange, not the spin vector \(\mathbf S\). Its action is
\begin{equation}
(S\psi)(x_{1},x_{2})=\psi(x_{2},x_{1}).
\end{equation}
A general function of two variables need not be symmetric, so this operation can produce a different wavefunction.

Swapping twice restores the original ordering:
\begin{equation}
S^{2}=I.
\end{equation}
The swap is reversible and preserves inner products, so it is represented by a unitary operator. If \(\psi\) is an eigenvector of \(S\) with eigenvalue \(\lambda\), then
\[
S^{2}\psi=\lambda^{2}\psi=\psi,
\]
and therefore
\begin{equation}
\lambda=\pm1.
\end{equation}

This conclusion identifies two exchange eigenspaces, but it does not say that every arbitrary two-particle wavefunction is already an eigenvector of \(S\). Any such wavefunction can be separated into its two parts,\footnote{This decomposition is an editorial algebraic clarification of the two eigenspaces implied by \(S^{2}=I\). The recorded argument states the two eigenvalues and then introduces the additional statistics principle.}
\begin{equation}
\psi_{\pm}=\frac12(I\pm S)\psi,
\end{equation}
which obey
\begin{equation}
S\psi_{\pm}=\pm\psi_{\pm}.
\end{equation}
Thus
\begin{align}
\psi_{+}(x_{2},x_{1})&=\psi_{+}(x_{1},x_{2}),\\
\psi_{-}(x_{2},x_{1})&=-\psi_{-}(x_{1},x_{2}).
\end{align}

Now comes an additional principle, not a consequence of \(S^{2}=I\) alone. For identical particles of any fixed species, the allowed states lie entirely in one of these exchange sectors: exchanging a pair always returns the same wavefunction, or always returns its negative. The \(+1\) sector is symmetric and the \(-1\) sector is antisymmetric. The physical interpretation of these two possibilities as the two particle statistics is the next step, and is left for the following discussion.
